% Direct-authority diplomatic French for printed p.171.
\SourcePage{171}

\LayerTarget{tohoku:proof:3.6.5:p2}
\noindent posantes irréductibles de $X$ [15, Chap. II, prop. 2]. Si $F$ est un
faisceau abélien sur $X$, on a un monomorphisme naturel de $F$ dans la somme
directe des $F_k=F_{X_k}$, d'où une suite exacte
\LayerTarget{tohoku:eq:irreducible-components-resolution}
\[
0\to F\to\oplus F_k\to R
\to0
\]
où $R$ est un faisceau dont le support est contenu dans
\LayerTarget{tohoku:def:singular-locus-union}
$Y=\displaystyle\bigcup\limits_{k\ne l}X_k\cap X_l$,
qui est de dimension combinatoire $\leqq n-1$. On en conclut une suite exacte
\LayerTarget{tohoku:eq:cohomology-component-sequence}
$H^{i-1}(X,R)\to H^i(X,F)\to\displaystyle\bigoplus\limits_k H^i(X,F_k)$.
Si $i>n$ d'oû $i-1>n-1$, on a
$H^{i-1}(X,R)=H^{i-1}(Y,R)=0$ d'après l'hypothèse de récurrence, il suffit
donc, pour prouver $H^i(X,F)=0$, de prouver que $H^i(X,F_k)=0$. Or on a
\LayerTarget{tohoku:eq:component-cohomology-identification}
$H^i(X,F_k)=H^i(X_k,F)$,
cela nous ramène à prouver le théorème pour l'espace irréductible $X_k$. Nous
supposerons donc $X$ irréductible. En vertu du corollaire à la proposition
3.6.3, il suffit de prouver que $H^i(X,\mathbf Z_U)=0$ pour $i>n$ et toute partie
ouverte $U$ de $X$. On peut supposer $U\ne\phi$, alors $Y=\complement U$ est
une partie fermée $\ne X$ de $X$, donc de dimension combinatoire
$\leqq n-1$. De la suite exacte
\LayerTarget{tohoku:eq:open-complement-sequence}
$0\to \mathbf Z_U\to \mathbf Z\to \mathbf Z_Y\to0$
on tire, puisque $\mathbf Z$ est flasque (cf. fin du 3.3),
\LayerTarget{tohoku:eq:open-support-cohomology-reduction}
$H^i(X,\mathbf Z_U)=H^{i-1}(X,\mathbf Z_Y)=H^{i-1}(Y,\mathbf Z)$,
qui est encore nul d'après l'hypothèse de récurrence.

\LayerTarget{tohoku:remark:3.6.5:1}
\indent\textsc{Remarques} 1. Le thème précédent généralise un théorème
antérieur de Serre [18].

\LayerTarget{tohoku:remark:3.6.5:2}
\indent 2. On peut trouver des espace de Zariski de dimension cohomologique
nulle, et de dimension combinatoire infinie, ou finie arbitrairement élevée.
\LayerTarget{tohoku:def:ordered-zariski-example}
Il suffit de considérer, sur un ensemble bien ordonné $X$ fini ou infini, la
topologie dont les fermés sont les ensembles de la forme $X_x$, où pour tout
$x\in X$, $X_x$ désigne l'ensemble des $y\in X$ tels que $y<x$.

\LayerTarget{tohoku:section:3.7}
\indent\textbf{3.7. La suite spectrale de Leray d'une application continue.} Soit
$f$ une application continue d'un espace $Y$ dans un espace $X$.
\LayerTarget{tohoku:def:relative-support-antifilter}
On suppose donné dans $Y$ un antifiltre $\Psi$ de partie fermées. Pour tout
ouvert $U\subset X$, soit $\Psi(U)$ l'antifiltre de parties fermées dans
$f^{-1}(U)$ formé des $A\subset f^{-1}(U)$ tels que, pour tout $x\in U$,
existe un voisinage $V\subset U$ de $x$ tel que l'adhérence de
$A\cap f^{-1}(V)$ soit dans $\Psi$. En particulier, $\Psi(X)$ désigne un
antifiltre de parties fermées de $Y$.
\LayerTarget{tohoku:def:relative-direct-image}
Soit $F$ un faisceau de groupes abéliens sur $Y$, considérons, pour tout ouvert
$U\subset X$, le groupe $\Gamma_{\Psi(U)}(f^{-1}(U),F)$. On vérifie facilement
que, pour les applications de restrictions évidentes (relatives à des
inclusion $V\subset U$), ces groupes forment un \textit{faisceau} sur $X$, noté
$f(F)$ et appelé \textit{image directe de $F$ par $f$ relativement à
$\Psi$}.
\LayerTarget{tohoku:def:ordinary-direct-image}
Si $\Psi$ est la famille de toutes les parties fermées de $Y$, on écrit
simplement $f_*(F)$ au lieu de $f_\Psi(F)$, et $f_*(F)$ est appelé
l'\textit{image directe du faisceau $F$ par $f$}. D'ailleurs, dans ce cas,
$\Psi(U)$ est l'ensemble de toutes les parties fermées de l'espace
$f^{-1}(U)$, donc
\LayerTarget{tohoku:eq:direct-image-sections}
$\Gamma(U,f_*(F))=\Gamma(f^{-1}(U),F)$
(définition qui a aussi sens si on suppose seulement que $F$ est un faisceau
d'ensembles). --- Par définition, on a dans le cas général :
\LayerTarget{tohoku:eq:relative-direct-image-sections}
\[
\Gamma(U,f(F))=\Gamma_{\Psi(U)}(f^{-1}(U),F).
\]

On vérifie aussitôt la formule :

\TohokuContinuation{172}
